\section{A positive elliptic kernel for \texorpdfstring{$x>1$}{x>1}}
\begin{theorem}\label{thm:positive}
For every real $x>1$, the sequence $(A_n(x))_{n\ge0}$ has a positive
compactly supported Stieltjes representing measure.
\end{theorem}
Fix $r=\sqrt{x}>1$ and set
\[
 h=1+8r^2-8r,\qquad H=1+8r^2+8r,\qquad c_r=\frac1{2\pi^2r}.
\]
Write $f_r$ for the density of the Legendre averaging law on $(-1,1)$.
We will continue $f_r$ to $\Omega=\C\setminus[1,\infty)$, express it as
the Cauchy transform of a positive density, and convert its Legendre
moments into power moments.

\subsection{The density and its continuation}
Throughout this section $\mathcal K$ uses a \emph{parameter}, not a modulus:
\[
 \mathcal K(w)=\int_0^{\pi/2}(1-w\sin^2\theta)^{-1/2}\,d\theta
             =R_F(0,1-w,1).
\]
Here $R_F$ is Carlson's symmetric integral, continued from
\[
 R_F(a,b,c)=\frac12\int_0^\infty
              \frac{dt}{\sqrt{(t+a)(t+b)(t+c)}}.
\]
Each factor uses its principal square root on the initial domain.
These conventions agree with \cite[Sections 19.16 and 19.25]{DLMF}
after replacing its squared modulus by our parameter.

For $-1<y<1$, put
\begin{equation}\label{eq:AB}
 d=\sqrt{(y+1)/2},\quad m=\frac{y-1}{4r},\quad
 A=(1-d)^2-m^2,\quad B=(1+d)^2-m^2 .
\end{equation}
Conditional on $c=\cos t$, the second angular variable has an arcsine
law. Its squared radicand factors as
\[
 16r^2(1+2rc+r^2)-(y-1-4r^2-4rc)^2
 =16r^2\bigl[d^2-(c-m)^2\bigr].
\]
The conditional law has no atoms, since its amplitude is at least
$4r(r-1)>0$. Integrating against the arcsine density of $c$ gives
\[
 f_r(y)=\frac1{4\pi^2r}\int_{m-d}^{m+d}
 \frac{dc}{\sqrt{(1-c^2)(c-m+d)(m+d-c)}} .
\]
Indeed, the roots of the conditional quadratic radicand are $m-d,m+d$,
both strictly inside $(-1,1)$ when $r>1$ and $-1<y<1$.
For example,
$1+m-d=(1-d)(1-(d+1)/(2r))>0$ and
$1-m-d=(1-d)(1+(d+1)/(2r))>0$.
The four-root substitution
\[
 t^2=\frac{(e_3-e_1)(c-e_2)}{(e_3-e_2)(c-e_1)},
 \qquad (e_1,e_2,e_3,e_4)=(-1,m-d,m+d,1),
\]
reduces this integral to the standard complete elliptic integral. Since
$(e_3-e_1)(e_4-e_2)=B$ and
$(e_3-e_2)(e_4-e_1)=4d$, it yields
\begin{equation}\label{eq:ellipticf}
 f_r(y)=c_r R_F(0,A,B)
       =c_r B^{-1/2}\mathcal K(4d/B).
\end{equation}

\begin{lemma}\label{lem:directanalytic}
The germ in \eqref{eq:ellipticf} extends to a single-valued holomorphic
function on $\Omega$.
\end{lemma}
\begin{proof}
On the two-sheeted surface $d^2=(z+1)/2$, the finite singularities of the
composition in \eqref{eq:ellipticf} can occur only above
$-1,1,h,H$. Indeed, $\mathcal K(p)=\pi\,{}_2F_1(1/2,1/2;1;p)/2$
continues along every path avoiding $p=0,1,\infty$, by the classical
hypergeometric equation \cite[15.10.1]{DLMF}. Here $p=4d/B$, and the
preimages are $d=0$, $A=0$, or $B=0$. The factorizations
\begin{equation}\label{eq:ABfactor}
 A=(d-1)^2\left(1-\frac{(d+1)^2}{4r^2}\right),\qquad
 B=(d+1)^2\left(1-\frac{(d-1)^2}{4r^2}\right)
\end{equation}
give precisely those projected points. Critical points of the parameter
map create no singularity in a composition, and the prefactor square root
continues along paths avoiding its zeros. Thus the selected germ continues
along every path in $\Omega\setminus\{-1\}$.

The apparent interior puncture is removable for this germ. The change
$c=m+d\cos\phi$ in the density integral gives, near $d=0$,
\[
 f_r=\frac1{4\pi^2r}\int_0^\pi
 \left[1-\left(\frac{d^2-1}{2r}+d\cos\phi\right)^2\right]^{-1/2}d\phi.
\]
At $d=0$ the expression inside the inner square has modulus
$1/(2r)<1$. For sufficiently small complex $d$ the integrand has a
single holomorphic square-root branch uniformly in $\phi$. Its integral
is holomorphic in $d$, and $\phi\mapsto\pi-\phi$ shows it is even.
It is therefore holomorphic in $z=2d^2-1$ and agrees with the initial germ
for real $d>0$. The generator loop about $-1$ has trivial monodromy.
Since $\Omega$ is simply connected and the fundamental group of
$\Omega\setminus\{-1\}$ is generated by that loop, the monodromy theorem
gives single-valuedness. The local integral then fills the puncture.
\end{proof}

\subsection{The boundary values}
The parameter versions of
the classical connection identities are
\begin{align}
 \mathcal K(z)&=(1-z)^{-1/2}\mathcal K\bigl(z/(z-1)\bigr),
                 && z<0,\label{eq:imaginaryK}\\
 \mathcal K(z)&=z^{-1/2}\bigl[
     \mathcal K(1/z)+i\mathcal K(1-1/z)\bigr],
                 && z\text{ above }(1,\infty).\label{eq:reciprocalK}
\end{align}
They hold as identities of the adjacent holomorphic germs; the lower-bank
version of \eqref{eq:reciprocalK} is its conjugate. See
\cite[19.7.2--19.7.3]{DLMF}, including its explicit upper/lower sign rule.
The convergent logarithmic expansion \cite[19.12.1]{DLMF} gives
\begin{equation}\label{eq:Klog}
 \mathcal K(1-\epsilon)
   =-\frac{\mathcal K(\epsilon)}{\pi}\log\epsilon+G(\epsilon),
\end{equation}
where $G$ is holomorphic near zero.

For real $y>1$, retain the positive square root in $d$.
The factorizations \eqref{eq:ABfactor} show that $A$ has a double zero
at $1$ and a simple zero at $h$, while
$B$ has a simple zero at $H$.
\begin{lemma}\label{lem:upperbank}
The upper boundary value of $f_r$ is given by the following holomorphic
germs adjacent to the indicated real intervals:
\begin{align}
 f_r(z)&=\frac{c_r}{\sqrt B}
   \bigl[\mathcal K(p)+2i\mathcal K(1-p)\bigr],
 &p=\frac{4d}{B},\quad&1<y<h,\label{eq:bankone}\\
 f_r(z)&=\frac{c_r}{\sqrt{4d}}
   \bigl[\mathcal K(q)+3i\mathcal K(1-q)\bigr],
 &q=\frac{B}{4d},\quad&h<y<H,\label{eq:banktwo}\\
 f_r(z)&=\frac{c_r}{\sqrt{-A}}
   \bigl[-2\mathcal K(t)+3i\mathcal K(1-t)\bigr],
 &t=\frac BA,\quad&H<y<\infty.\label{eq:bankthree}
\end{align}
The parameters $p,q,t$ respectively lie in $(0,1)$ on their displayed
real intervals, and all displayed real prefactors use positive roots.
\end{lemma}
\begin{proof}
Near $1$,
\[
 A=\frac{1-r^{-2}}{16}(z-1)^2+O((z-1)^3),\qquad B=4+O(z-1).
\]
Going around $1$ through the upper half-plane makes $A/B$ wind clockwise
once about zero. Equation~\eqref{eq:Klog} adds exactly
$2i\mathcal K(A/B)$, proving \eqref{eq:bankone} near $1$.
On compact subintervals of $(1,h)$, both $p$ and $1-p$ stay in $(0,1)$.
Overlapping neighborhoods therefore propagate the same germ without an
additional period.

At $h$,
\[
 A'(h)=-\frac{(r-1)^2}{r(2r-1)}<0,\qquad
 p'(h)=-A'(h)/B(h)>0 .
\]
An upper semicircle about $h$ thus sends $p$ through the upper side of $1$.
Apply \eqref{eq:reciprocalK} to $\mathcal K(p)$ and
\eqref{eq:imaginaryK} to $\mathcal K(1-p)$:
\[
 \mathcal K(1-p)=p^{-1/2}\mathcal K(1-1/p).
\]
Substitution in \eqref{eq:bankone} gives \eqref{eq:banktwo}, including
the coefficient $3$. The square-root factors agree on these germs because
$Bp=4d$ and they agree before crossing. Again, compact subintervals of
$(h,H)$ keep $q,1-q$ strictly between $0$ and $1$.

At $H$,
\[
 B'(H)=-\frac{(r+1)^2}{r(2r+1)}<0,\qquad q'(H)=B'(H)/(4d(H))<0.
\]
Hence $q$ crosses zero through the lower half-plane and $1-q$ crosses $1$
through the upper half-plane. Apply \eqref{eq:imaginaryK} to $\mathcal K(q)$
and \eqref{eq:reciprocalK} to $\mathcal K(1-q)$. With
\[
 \frac{q}{q-1}=t,\qquad \frac1{1-q}=1-t,\qquad
 (4d)(1-q)=-A,
\]
the real coefficient becomes $1-3=-2$, giving \eqref{eq:bankthree}.
The parameters $t,1-t$ remain in $(0,1)$ to the right of $H$.
Overlapping neighborhoods propagate this germ to every finite point of
that bank. These are valid continuations of the selected germ in
$\Omega$, where Lemma~\ref{lem:directanalytic} supplies single-valuedness.
\end{proof}

Define $\rho_r(y)=\operatorname{Im}f_r(y+i0)/\pi$.
Lemma~\ref{lem:upperbank} gives the explicitly positive density
\begin{equation}\label{eq:rho}
 \rho_r(y)=\frac1{2\pi^3r}
 \begin{cases}
 2R_F(0,4d,B),&1<y<h,\\
 3R_F(0,4d,B),&h<y<H,\\
 3R_F(0,-A,-B),&H<y<\infty.
 \end{cases}
\end{equation}
Thus $\rho_r$ has a finite right limit at $1$, a finite jump at $h$,
and at most a logarithmic singularity at $H$. At infinity,
$-A,-B\sim y^2/(16r^2)$, so homogeneity gives $\rho_r(y)=O(1/y)$.

\subsection{The Cauchy representation}
\begin{lemma}\label{lem:dispersion}
For every $z\in\Omega$,
\begin{equation}\label{eq:dispersion}
 f_r(z)=\int_1^\infty\frac{\rho_r(a)}{a-z}\,da .
\end{equation}
\end{lemma}
\begin{proof}
The integral converges by the local logarithmic bounds and
$\rho_r(a)=O(1/a)$. To apply Cauchy's theorem, we need a bound that is
uniform around a large circle, including angles close to the cut.

Take $R_0=1000r^2$ and the connected upper exterior domain
$D=\{z:\operatorname{Im}z>0,\ |z|>R_0\}$. Choose
$d=\sqrt{(z+1)/2}$ by continuation from the positive bank. The expansion
\[
 A=-\frac{z^2}{16r^2}
   +\left(\frac12+\frac1{8r^2}\right)z
   +\frac32-\frac1{16r^2}-2\sqrt{(z+1)/2}
\]
gives $|A+z^2/(16r^2)|\le5R/4$, $R=|z|>R_0$, at every angle. Thus
\[
 \frac{R^2}{32r^2}\le |A|\le\frac{3R^2}{32r^2},\qquad
 \frac{\sqrt R}{2}\le |d|\le\sqrt R .
\]
Put $w=1-t=-4d/A$ and $A=-z^2(1+e)/(16r^2)$. Then
\[
 \frac{64r^2}{3R^{3/2}}\le |w|
       \le\frac{128r^2}{R^{3/2}}<\frac12,\qquad |e|\le\frac1{50}.
\]
Use $0<\arg z<\pi$ and the small analytic logarithms of $1+1/z$ and
$1+e$. The logarithm of $w$ continued from the positive bank is
\begin{equation}\label{eq:continuedlog}
 \operatorname{Log}w=\log\frac{64r^2}{\sqrt2}
 -\frac32\operatorname{Log}z+\frac12\log(1+1/z)-\log(1+e).
\end{equation}
In particular,
\[
 |\operatorname{Im}\operatorname{Log}w|
 \le\frac{3\pi}{2}+\tfrac12\arcsin(1/1000)+\arcsin(1/50)<5,
 \qquad |\operatorname{Log}w|\le\tfrac32\log R+5.
\]
Here $R^{-3/2}<|w|<1$ controls the real part. This logarithm is
\emph{continued}, not replaced by the principal logarithm: its imaginary
part can approach $-3\pi/2$ near the negative real axis.
The square root of $-A$ on $D$ is chosen positive at the positive bank;
its modulus is $|A|^{1/2}$.

In \eqref{eq:Klog}, the convergent series \cite[19.12.1--19.12.3]{DLMF} is
\[
 G(w)=\sum_{j\ge0}\left(\frac{(1/2)_j}{j!}\right)^2
       [\psi(j+1)-\psi(j+1/2)]w^j.
\]
Its coefficients satisfy $((1/2)_j/j!)^2\le1$ and
$0<\psi(j+1)-\psi(j+1/2)\le2\log2$. Hence for $|w|\le1/2$,
\[
 |\mathcal K(w)|\le\pi/\sqrt2,\qquad |G(w)|\le4\log2.
\]
Equation~\eqref{eq:bankthree} becomes
\begin{equation}\label{eq:fullcircle}
 f_r(z)=\frac{c_r}{\sqrt{-A}}
 \left[\frac2\pi\mathcal K(w)\operatorname{Log}w
       -2G(w)+3i\mathcal K(w)\right].
\end{equation}
All terms on the right are holomorphic on $D$ with the specified branches.
The identity theorem extends this identity from its positive-bank
neighborhood to all $D$. It gives the explicit full-semicircle bound
\[
 |f_r(z)|\le\frac{\sqrt{32}}{2\pi^2R}
 \left[\sqrt2\left(\tfrac32\log R+5\right)
             +8\log2+\frac{3\pi}{\sqrt2}\right].
\]
Schwarz reflection gives the lower semicircle, and holomorphy supplies the
negative-axis limit. Thus the estimate is uniform on the entire slit
circle, and the large-circle Cauchy integral vanishes.

At $1$, $A/B$ has a double zero and \eqref{eq:Klog} gives
$f_r(z)=O(1+|\log|z-1||)$. At $h$, $1-p$ has a simple zero in
\eqref{eq:bankone}; at $H$, $q$ has a simple zero in
\eqref{eq:banktwo}. The same logarithmic expansion therefore gives
$f_r(z)=O(1+|\log|z-a||)$ near each $a\in\{1,h,H\}$ on either
side of the cut. An arc of radius $\epsilon$ contributes
$O(\epsilon(1+|\log\epsilon|))$, which tends to zero.
Cauchy's theorem
on the slit plane now gives \eqref{eq:dispersion}, with
$(f_+-f_-)/(2\pi i)=\rho_r$ and denominator $a-z$.
\end{proof}

\subsection{The positive remainder measure}
Common's positive Froissart--Gribov theorem
\cite[CERN preprint, Section 4, Theorem 5]{Common} converts positive
$Q_n$ mixtures into Stieltjes moments. We give the integral calculation
for the present density, whose spectral support reaches $1$ and whose
total spectral mass need not be finite.

Use the standard normalization
\[
 Q_n(a)=\frac12\int_{-1}^1\frac{P_n(y)}{a-y}\,dy,\qquad a>1.
\]
Heine's integral is
\[
 Q_n(a)=\int_0^\infty
       \bigl(a+\sqrt{a^2-1}\cosh u\bigr)^{-n-1}\,du.
\]
Set $t=(a+\sqrt{a^2-1}\cosh u)^{-1}$. Since
$du=-dt/(t\sqrt{1-2at+t^2})$, this gives
\begin{equation}\label{eq:Qpositive}
 Q_n(a)=\int_0^{a-\sqrt{a^2-1}}
       \frac{t^n}{\sqrt{1-2at+t^2}}\,dt .
\end{equation}
In \cite[14.12.13--14.12.14]{DLMF}, Olver's normalized function is
$Q_n/n!$; multiplying both formulas by $n!$ yields precisely this
normalization.

\begin{proof}[Proof of Theorem~\ref{thm:positive}]
First take $n=0$. All kernels in \eqref{eq:dispersion} and
\eqref{eq:Qpositive} are nonnegative on their real integration domains,
so Tonelli's theorem gives
\[
 \int_{-1}^1 f_r(y)\,dy
 =2\int_1^\infty\rho_r(a)Q_0(a)\,da
 =\int_0^1 w_r(t)\,dt<\infty,
\]
where $w_r$ is defined below. The first integral is the probability
$\mu_r([-1,1))$; the endpoints have zero mass by the conditional
arcsine representation. For general $n$, put
$b_n=\int_{-1}^1P_n(y)f_r(y)\,dy$. The bounds $|P_n(y)|\le1$ and
$0\le t^n\le1$ reduce absolute integrability to the case $n=0$.
Fubini's theorem now gives
\begin{align}
 b_n&=2\int_1^\infty \rho_r(a)Q_n(a)\,da
      =\int_0^1 t^n w_r(t)\,dt,\label{eq:positiveremainder}\\
 w_r(t)&=2\int_1^{(t+t^{-1})/2}
       \frac{\rho_r(a)}{\sqrt{1-2at+t^2}}\,da\ge0,
            \qquad 0<t<1.\label{eq:wmeasure}
\end{align}
The domain in \eqref{eq:wmeasure} follows from
$t\le a-\sqrt{a^2-1}$ if and only if
$1\le a\le(t+t^{-1})/2$. Thus $w_r(t)\,dt$ is a finite positive
measure on $[0,1]$ with moments $b_n$.

The contribution from $Y_r\ge1$ already has the positive Laplace-mixture
measure $\nu_+$, supported on
$[H-\sqrt{H^2-1},H+\sqrt{H^2-1}]$.
Adding the positive measure in \eqref{eq:wmeasure} proves that
$\nu_++w_r(t)\,dt$ is a positive compact Stieltjes measure with moments
$A_n(x)$.
\end{proof}

\begin{proof}[Proof of Theorem~\ref{thm:classification}]
Theorem~\ref{thm:main} excludes $0<x<1$ and
Theorem~\ref{thm:positive} covers $x>1$. Edgar's prior theorem
\cite{Edgar} covers $x=1$. At $x=0$, $A_n(0)=1$ is represented by
the point mass at $1$.
\end{proof}
